\section{Introduction}\label{sec:intro}

Let $\A$ be a finite alphabet and $\sh$ the left shift on $\A^{\Z}$. By the Curtis--Hedlund--Lyndon theorem \cite{Hedlund1969} the endomorphisms of the full shift are the block maps, and the automorphisms are the invertible block maps. The group $\Aut(\A^{\Z})$ is enormous: it contains every finite group and free groups \cite{BLR1988}, and every countable, locally finite, residually finite group \cite{KimRoush1990}, and its center consists of the powers of the shift \cite{Ryan1972}. A reversible cellular automaton $F\in\Aut(\A^{\Z})$ determines a two-dimensional subshift, its \emph{spacetime} $X_F\subset\A^{\Z^2}$, whose rows are the configurations $F^n(c)$, $n\in\Z$ \cite{CFK2019}. Its automorphism group is the centralizer of $F$ in $\Aut(\A^{\Z})$ (Lemma~\ref{lem:spacetime}), and it always contains $\langle\sh,F\rangle$. This paper is about automata for which it is much larger.

The automata we consider are the second-order elementary rules, built by the second-order technique of Fredkin \cite{TM1990,Wolfram2002}; as a family they were introduced and studied systematically by Takesue \cite{Takesue1987,Takesue1989}, under the name of elementary reversible cellular automata. For an elementary rule $f\colon\{0,1\}^3\to\{0,1\}$ define, on the alphabet $\A=\{0,1\}^2$ with cells written $(x,y)$,
\begin{equation}\label{eq:F}
F_f(c)_i=\bigl(y_i,\;x_i\oplus f(y_{i-1},y_i,y_{i+1})\bigr).
\end{equation}
Reading $y$ as the current row and $x$ as the previous one, this is the recursion $y_{t+1}=y_{t-1}\oplus f(y_t)$, which is reversible for every $f$, with $F_f^{-1}(c)_i=\bigl(y_i\oplus f(x_{i-1},x_i,x_{i+1}),\,x_i\bigr)$. Following Takesue \cite{Takesue1987}, one writes $f$R for this automaton; $37$R is Wolfram's standard example of a reversible automaton with complex behaviour \cite[pp.~440 and 452--457]{Wolfram2002}. Under reflection and the cellwise complement the $256$ rules fall into $88$ conjugacy classes, the classes of Takesue \cite[Table~I]{Takesue1989}, and $\Aut(X_{F_f})$ depends only on the class (Section~\ref{sec:prelim}).

For some of these automata the automorphism group is as small as possible. For rule $30$, for all non-affine rules permutive in an outer variable, and for the classes of $63$ and $187$, $\Aut(X_{F_f})=\langle\sh,F_f\rangle\cong\Z^2$ \cite{PaperB}; the same holds for the classes of $110$, $14$, $46$ and $79$ \cite{PaperA}. The classification of all $88$ classes is the subject of \cite{PaperD}. Here we study the opposite behaviour. The spacetime $X_F$ has zero $\Z^2$-entropy, so Hochman's theorem that positive entropy forces every finite group into the automorphism group of a $\Z^2$ shift of finite type \cite{Hochman2010} never applies to it, and something else must supply the room. In every case we find, it is a local structure.

For most of the classes treated here the local structure is a \emph{wall} in the sense of Takesue \cite[Section~4]{Takesue1989}: a pattern, fixed or periodic in time, that appears at a fixed position irrespective of its surroundings. Takesue relates walls to local conservation laws, exhibits walls in $40$ of the $88$ classes (the $38$ classes of \cite[Fig.~2, Table~III]{Takesue1989} and those of $27$ and $59$), and observes that the two sides of a wall cannot communicate, so that a large system with walls is divided into parts that do not interact and is not mixing. This observation precedes the formalism of blocking words and equicontinuity points \cite{Kurka1997,BT2000}, of which walls are instances (Remark~\ref{hw:kurka}). The walls of the second-order elementary automata are therefore known, and we refer to \cite{Takesue1989} for them. What this paper adds is their consequence for the automorphism group: the embeddings of $\bigoplus_mS_m$ by permutations along trains of walls; a proof that the list of fixed walls is complete for words of every length; a bit hidden in words built from the periodic walls, with the flips that it gives; the census of rule $59$R at radius $2$; and the results on rule $54$. Where there are no walls, it adds the results on the linear and affine rules.

\subsection*{Fixed walls}
The basic mechanism is a word that the automaton fixes in every context.

\begin{theorem}\label{thm:walls}
Let $F$ be a reversible cellular automaton of radius $1$ on $\A^{\Z}$ and $W$ a nonempty word such that $F(c)$ contains $W$ at every position at which $c$ contains $W$. Let $C\subseteq\A$ be the set of symbols $s$ for which $W\,s\,W$ contains $W$ only at its two ends. If $|C|\ge2$, then $\Aut(X_F)$ contains a copy of $\bigoplus_{m\ge1}S_m$, hence of every finite group.
\end{theorem}

The copies of $W$ act as walls. A symbol between two walls evolves by a fixed permutation of $C$, whatever lies beyond the walls, so the contents of a run of such cells can be permuted without disturbing the dynamics. For the second-order elementary automata such words are known: the shortest fixed words of the $31$ classes of Table~\ref{tab:walls}, and the word $0330$ for the class of $133$, are the fixed walls (i)--(ix) of Takesue \cite[Fig.~2, Table~III and p.~391]{Takesue1989}, who examined the words of length at most $6$. We verify the hypothesis $|C|\ge2$ for them, and we prove that no other class has a fixed word of any length (Proposition~\ref{cl:prop:133}).

\begin{corollary}\label{cor:48}
For the $81$ elementary rules $f$ in the $31$ classes listed in Table~\ref{tab:walls}, $\Aut(X_{F_f})$ contains every finite group. These include $54$, $37$, $22$ and $73$. For the class $\{133,161\}$ the same holds, through the fixed word $0330$ of length $4$. No other elementary rule has a word of any length fixed by $F_f$ in every context.
\end{corollary}

\subsection*{Rule 54}
Rule $54$ is, with rule $110$, one of the two elementary rules for which Wolfram reports class~4 behaviour, rule $110$ as the clear example and rule $54$ as a possible one \cite[pp.~693--697]{Wolfram2002}. The two second-order automata behave in opposite ways: $\Aut(X_{F_{110}})=\langle\sh,F_{110}\rangle$ \cite{PaperA}, while for $54$R the pattern $33$ (writing $3$ for the cell $(1,1)$) is a wall, the wall (iv) of \cite[Table~III]{Takesue1989}, and the walls give the following.

\begin{theorem}\label{thm:finite}
For every $m\ge1$ there is an injective homomorphism $\Psi_m\colon S_m\to\Aut(X_{F_{54}})$ whose image consists of block maps of radius at most $3m+3$, and the images for different $m$ commute, so $\bigoplus_{m\ge1}S_m$ embeds in $\Aut(X_{F_{54}})$.
\end{theorem}

The relabellings of the segments between walls are automorphisms of smaller radius: a permutation of the contents of the segments of a given length $L$ is an automorphism whenever it commutes with the permutation that $F$ induces on these contents, and these maps give an embedding of the direct sum over $L$ of the centralizers of the induced permutations (Theorem~\ref{thm:54}). The smallest of them is $D$, of radius $2$ and order $3$, which acts as the $3$-cycle $0\to1\to2\to0$ on every cell that is a single-cell segment between two walls and as the identity elsewhere. At radius $2$ an exhaustive computation shows that the walls account for everything: the automorphisms of radius at most $2$ are exactly the $13$ maps $\sh^aF^b$ with $|a|+|b|\le2$ and the $10$ maps $F^bD^k$ with $|b|\le 2$, $k\in\{1,2\}$ (Theorem~\ref{thm:54census}). At radius $3$ this is no longer so: there is an involution $\tau$ of radius $3$ that commutes with $F$, moves the boundary between a long wall and a single cell next to it, and does not commute with $D$ (Theorem~\ref{thm:tau}). A partial enumeration of the automorphisms of radius $3$ (Section~\ref{sec:54r3}) finds that every automorphism it produces differs from some $\sh^aF^b$ only on windows that contain a wall, and that most of them, like $\tau$, do not preserve walls.

\subsection*{Walls that are not fixed}
A wall need not be fixed in every context. Takesue's list contains walls that are periodic in time and not fixed, in the classes of $27$, $59$, $179$, $127$, $223$, $126$ and $219$ \cite[Section~4, (4.1)--(4.3), Fig.~2 and Table~III, walls (xii), (xiii), (xv), (xvi)]{Takesue1989}. None of these seven classes has a word fixed in every context, so Theorem~\ref{thm:walls} does not apply to them. In Section~\ref{sec:hidden} we give a general mechanism (Theorem~\ref{hw:general}) for walls whose image never depends on the context, except at cells that do not matter, and which carry one bit that $F$ cannot see. Permuting these hidden bits along chains of walls gives every finite group again. The mechanism applies to words built from Takesue's periodic walls in the classes of $59$, $27$, $127$, $179$ and $223$ (Theorems~\ref{hw:59} and~\ref{hw:27}, Proposition~\ref{hw:tp}). For $219$R and $126$R the walls rotate through the alphabet in every context, and permuting the cells between them again embeds every finite group (Theorem~\ref{hw:219}, Proposition~\ref{hw:126}).

\subsection*{The affine rules 90 and 165}
Rule $90$ has no blocking word \cite{Kurka1997,BT2000}, hence no wall of any of the kinds above (Corollary~\ref{lin:noblock}), and the automorphisms of $X_{F_{90}}$ of radius at most $2$ are all affine (Section~\ref{sec:linear}). Nevertheless $F_{90}$ has two factor maps onto the binary full shift that carry $F$ to $\sh^{\pm1}$, with matching embeddings (Theorem~\ref{thm:90}), and these carry marker automorphisms of the binary full shift into $X_{F_{90}}$: by Theorem~\ref{lin:movers}, $(\bigoplus_mS_m)^2$ embeds in $\Aut(X_{F_{90}})$ and in $\Aut(X_{F_{165}})$. Theorem~\ref{thm:90} also describes the block maps of radius at most $3$ that commute with $F_{90}$, invertible or not: they are the sums of an affine map and of maps that factor through the two factor maps.

\subsection*{The other affine rules}
For the affine rules whose characteristic polynomial $t^2+pt+1$ is irreducible and whose spacetime action is mixing, the isomorphism rigidity theorem of Kitchens and Schmidt \cite{KS2000} forces every automorphism to be affine (Lemma~\ref{aff:lem:R}), and the automorphism group becomes a unit group, which we compute. Write $\kappa$ for the cellwise complement.

\begin{theorem}\label{thm:affine}
Every automorphism of $X_{F_f}$ is affine for the rules $f$ of the four classes below, and:
\begin{enumerate}
\item for $f\in\{170,240\}$, $\Aut(X_{F_f})=\langle F,h\rangle\cong\Z^2$ with $h=\mathrm{id}+F$, and $h^2=\sh^{\pm1}F$, so $\langle\sh,F\rangle$ has index $2$; for $f\in\{15,85\}$ the same holds with $h(c)=c+F(c)+(1,1)^{\Z}$;
\item for $f\in\{60,102\}$, $\Aut(X_{F_f})=\langle\sh,F\rangle\times\langle\kappa\rangle\cong\Z^2\times\Z/2$;
\item for $f\in\{153,195\}$, $\Aut(X_{F_f})=\langle\sh,F\rangle\times\langle\kappa\rangle\cong\Z^2\times\Z/2$;
\item for $f=150$, $\Aut(X_{F_f})=\langle\sh,F,h\rangle\cong\Z^3$ with $h=\mathrm{id}+\sh F$; for $f=105$ the same holds with $h(c)=c+\sh F(c)+(1,1)^{\Z}$.
\end{enumerate}
\end{theorem}

For $170$ the spacetime, with the action of $F$ and $\mathrm{id}+F$, is conjugate to Ledrappier's three-dot shift, with $\langle\sh,F\rangle$ acting as a subgroup of index $2$ of its shift action (Remark~\ref{aff:rem:ledrappier}). The hypotheses fail exactly for the other affine classes, those of $0$, $51$, $90$, $165$ and $255$, and each of them has non-affine automorphisms (Remark~\ref{aff:rem:consist}). For $255$, $F$ adds $1$ modulo $4$ in every cell, and $\Aut(X_{F_{255}})\cong\Z/4\times\mathcal G_4$, where $\mathcal G_4$ is a group of automorphisms of the full shift on $\Z/4$ that contains $\bigoplus_mS_m$ (Proposition~\ref{aff:prop:255}).

\subsection*{Summary}
Collecting these results, the following $42$ classes of second-order elementary automata, comprising $102$ rules, have the property that $\bigoplus_{m\ge1}S_m$ embeds in $\Aut(X_{F_f})$, so that $\Aut(X_{F_f})$ contains every finite group and is not virtually $\Z^n$ for any $n$:
\begin{itemize}
\item the $31$ classes ($81$ rules) of Table~\ref{tab:walls}, with a word of length at most $3$ fixed in every context, and the class of $133$ ($2$ rules), with one of length $4$ (Theorem~\ref{thm:walls}, Corollary~\ref{cor:48});
\item the classes of $59$ and $27$ ($8$ rules), with hidden-bit walls (Theorems~\ref{hw:59} and~\ref{hw:27});
\item the classes of $127$, $179$ and $223$ ($6$ rules), with time-periodic walls (Proposition~\ref{hw:tp});
\item the classes of $219$ and $126$ ($2$ rules), with rotating walls (Theorem~\ref{hw:219}, Proposition~\ref{hw:126});
\item the affine classes of $90$ and $165$ ($2$ rules), where $(\bigoplus_mS_m)^2$ embeds (Theorem~\ref{lin:movers});
\item the affine class of $255$ ($1$ rule), through marker permutations of the difference sequence (Proposition~\ref{aff:prop:255}).
\end{itemize}
The first four items comprise $39$ classes. Their walls are those of Takesue or are built from them, and with the class of $255$, which also has walls in his list, they are the $40$ classes in which he exhibits walls \cite[Section~4]{Takesue1989}.
Apart from the description $\Aut(X_{F_{255}})\cong\Z/4\times\mathcal G_4$, we do not determine the whole automorphism group in these cases. For the other $46$ classes, the status of each is recorded in \cite{PaperD}: for the class of $95$ the automorphism group is determined in \cite{PaperE}, and it also contains $\bigoplus_{m\ge1}S_m$, through automorphisms that permute moving particles and not through a wall; $11$ classes have $\Aut(X_{F_f})=\langle\sh,F_f\rangle\cong\Z^2$ \cite{PaperA,PaperB,PaperE}, the four affine classes of Theorem~\ref{thm:affine} have the groups given there, two further classes have the cellwise symmetry $\kappa$, and for $28$ classes an exhaustive search finds no automorphism outside $\langle\sh,F_f\rangle$ of radius at most $3$.

\subsection*{Blocking words and linearity}
In the nonlinear examples of this paper the automorphisms outside $\langle\sh,F\rangle$ all come from blocking words with a degree of freedom next to the blocking column (Remark~\ref{hw:kurka}). The following remark records what the affine rules say about the converse.

\begin{remark}\label{rem:dichotomy}
Without a hypothesis on $f$, the absence of a blocking word does not force $\Aut(X_{F_f})$ to be virtually $\Z^2$. The affine rules $90$ and $165$ have no blocking word (Corollary~\ref{lin:noblock}), yet $(\bigoplus_mS_m)^2$ embeds in their automorphism groups (Theorem~\ref{lin:movers}). Rule $150$ is a counterexample of a different kind. It has no blocking word, and its commuting block maps of radius at most $3$ are all linear (Proposition~\ref{prop:150}); but $\mathrm{id}+\sh F_{150}$ is an automorphism outside $\langle\sh,F_{150}\rangle$, and $\langle\sh,F_{150},\mathrm{id}+\sh F_{150}\rangle\cong\Z^3$ (Proposition~\ref{prop:150Z3}), so $\Aut(X_{F_{150}})$ is not virtually $\Z^2$; in fact $\Aut(X_{F_{150}})=\langle\sh,F_{150},\mathrm{id}+\sh F_{150}\rangle$ (Theorem~\ref{thm:affine}). In all three the room comes from linearity, not from a local structure.
\end{remark}

Nor does the absence of a blocking word force $\Aut(X_{F_f})$ to be virtually $\Z^2$ when $f$ is not affine: the rule $95$ is not affine, $F_{95}$ has no blocking word, and $\Aut(X_{F_{95}})$ contains every finite group \cite{PaperE}. There the room comes from automorphisms that permute particles moving at light speed, and neither from a local structure nor from linearity. We do not know another example of this kind among the elementary rules.

\begin{question}\label{q:dichotomy}
Let $f$ be an elementary rule that is not affine and does not belong to the class of $95$. If the second-order automaton $F_f$ has no blocking word \cite{Kurka1997,BT2000}, is $\Aut(X_{F_{f}})$ virtually $\Z^2$?
\end{question}

Having no blocking word is stronger than having no word fixed in every context; the rules of Section~\ref{sec:hidden} show that the difference matters (Remark~\ref{hw:kurka}). For $110$R and $30$R, whose automorphism groups are $\Z^2$ \cite{PaperA,PaperB}, a single changed cell produces a difference that spreads in both directions at linear speed, whereas in $54$R it spreads only a few cells (Table~\ref{tab:spread}).

\subsection*{Related work}
\emph{Large automorphism groups.} The automorphism group of a mixing one-dimensional shift of finite type contains every finite group \cite{BLR1988}. In two dimensions the same holds under strong irreducibility \cite{Ward1994} and under positive entropy \cite{Hochman2010}. Hochman also notes that zero entropy does not prevent a large automorphism group, with the example of a shift whose columns are constant; that shift is the spacetime of the identity automaton. Theorems~\ref{thm:walls} and~\ref{thm:finite} are of the same kind, but for nonlinear automata of infinite order, where the room comes from a conserved local structure. For linear automata the algebraic theory of Kitchens and Schmidt applies \cite{KS1989}; their isomorphism rigidity theorem \cite{KS2000} is what determines the automorphism groups in Theorem~\ref{thm:affine}, and for Ledrappier's shift \cite{Ledrappier1978} the automorphism group is $\Z^2$ \cite{Ward1994}, \cite[Theorem~7]{BRY2018}, \cite[Example~8.20]{Schmidt2002}.

\emph{Centralizers and blocking words.} Ryan's theorem identifies the centralizer of the whole automorphism group of a full shift with the powers of the shift \cite{Ryan1972}, and finitary versions find finite sets of automorphisms with the same centralizer \cite{Salo2019Ryan,Kopra2022}; a spacetime automorphism group is the centralizer of a single automaton, which these results do not cover. The term blocking word is due to Blanchard and Tisseur \cite{BT2000}; the relation of blocking words to equicontinuity is in the proof of \cite[Theorem~3]{Kurka1997}. For the second-order elementary automata, that a wall separates the two sides of a configuration was observed by Takesue \cite{Takesue1987}, \cite[Section~4]{Takesue1989}, before these notions were introduced.

\emph{Rule 54.} The particles and collisions of the elementary rule $54$ were catalogued by Boccara, Nasser and Roger \cite{BNR1991} and by Mart\'inez, Adamatzky, Chen and Chua \cite{MAM2013}. A different reversible automaton, introduced by Bobenko, Bordemann, Gunn and Pinkall \cite{BBGP1993} and studied extensively as an integrable model of statistical mechanics \cite{BKP2021}, is also called Rule~54 because it uses the same local function $b\oplus(a\vee c)$, but it updates the two sublattices alternately in place. In the present notation it is one of the two sublattices of a spacetime diagram of $250$R \cite{BKP2021}, which is the conjugate of $95$R under the cellwise complement: it is the half system of the class of $95$, whose automorphism group is determined in \cite{PaperE}. It is not the second-order automaton $F_{54}$ studied here, which has the extra term from the previous row; its solitons move at speed $\pm1$, while in $F_{54}$ the structures that carry the automorphisms are the stationary walls.

\subsection*{Organisation}
Section~\ref{sec:prelim} fixes notation, proves Lemma~\ref{lem:spacetime}, and shows that $\langle\sh,F_f\rangle\cong\Z^2$ unless $f$ reads only its middle argument. Section~\ref{sec:method} describes the SAT encoding of the commutation condition, the invertibility test and the verification. Section~\ref{sec:54} treats rule $54$: it proves Theorem~\ref{thm:finite}, and it states and proves Theorems~\ref{thm:54} and~\ref{thm:tau}. Section~\ref{sec:walls} proves Theorem~\ref{thm:walls} and Corollary~\ref{cor:48}, including the class of $133$. Section~\ref{sec:hidden} treats walls with hidden bits, rules $59$ and $27$, the time-periodic walls of rules $127$, $179$ and $223$, the rotating walls of rules $219$ and $126$, and blocking words. Section~\ref{sec:linear} treats the linear rules $90$ and $150$ (Theorem~\ref{thm:90}, Proposition~\ref{prop:150}) and proves that $(\bigoplus_mS_m)^2$ embeds in the automorphism groups of $90$R and $165$R. Section~\ref{sec:affine} proves Theorem~\ref{thm:affine} and treats rule $255$. Section~\ref{sec:questions} lists open questions. The appendix describes the computations.
